\begin[standout]{frame}
    Quelques notions sur les systèmes distribués
\end{frame}

\begin{frame}{Architecture distribuée : les bases}

    Contenu :

    Définir simplement : nœuds, réplication, partitionnement, tolérance aux pannes.
    Schéma d’une architecture distribuée (ex : master-slave vs peer-to-peer).
    Exemple concret : Comment une requête est traitée dans un cluster MongoDB.

    Format : Un schéma SVG ou un diagramme Mermaid.
\end{frame}

\begin{frame}{Systèmes distribués}
    \begin{definitionbox}
        Un système distribué est un ensemble d'ordinateurs autonomes qui apparaissent à l'utilisateur comme un système unique.
    \end{definitionbox}

    \pause

    \begin{itemize}
        \item Les systèmes distribués sont utilisés pour stocker et traiter de grandes quantités de données.
        \item Ils permettent de répartir la charge de travail sur plusieurs machines.
        \item Ils permettent d'assurer la disponibilité et la tolérance aux pannes.
    \end{itemize}
\end{frame}
